\documentclass[aspectratio=169,10pt]{beamer}
\usetheme{DSI}

\usepackage{listings}
\lstdefinestyle{slidesql}{
  language=SQL,
  morekeywords={WITH,RECURSIVE},
  basicstyle=\ttfamily\small\color{UHblack},
  keywordstyle=\color{red1}\bfseries,
  commentstyle=\color{UHblack!55}\itshape,
  stringstyle=\color{red2},
  identifierstyle=\color{UHblack},
  columns=fullflexible,
  keepspaces=true,
  showstringspaces=false
}



%%If you want the change the logo to your research group logo, follow the steps in line 168 of beamerthemeDSI.sty. If you want no logo you can just comment lines 164-166

\title{Analyzing loop hoisting optimizations for recursive SQL queries}
\author{Olaf Van Bylen}
\newcommand{\collaboration}{In collaboration with Frank Neven, Brecht Vandevoort and Stijn Vansummeren}
% \institute{Data Science Institute - Hasselt University}
\date{December 4, 2026}

% Card spacing can be adjusted locally without losing automatic text sizing.
\newcommand{\motivationtitlefont}{\large\bfseries}
\newcommand{\motivationbodyfont}{\small}
\newcommand{\motivationpadding}{10pt}
\newcommand{\motivationgap}{8pt}
\newcommand{\motivationtextgap}{4pt}

% Cards size themselves to the wrapped title and body, with consistent padding.
\newcommand{\motivationcard}[3]{%
  \par\noindent
  \begin{tikzpicture}
    \node[fill=#3,draw=UHblack!9,rounded corners=4pt,
      text width=\dimexpr\linewidth-0.8cm-1pt\relax,
      inner xsep=0.4cm,inner ysep=\motivationpadding,outer sep=0pt,
      align=left,font=\motivationbodyfont] (card) {%
      {\color{UHblack}\motivationtitlefont #1\par}
      \vspace{\motivationtextgap}
      {\color{UHblack!75}#2\par}%
    };
    \draw[UHred,line width=2pt,line cap=round]
      ([xshift=1pt,yshift=-8pt]card.north west) --
      ([xshift=1pt,yshift=8pt]card.south west);
  \end{tikzpicture}%
  \par\addvspace{\motivationgap}
}

\begin{document}
% --------------------------------------------------- Slide --
\begin{frame}
  \titlepage
\end{frame}

\begin{frame}{Contents}
 \tableofcontents
\end{frame}

\begin{frame}{Motivation}
\renewcommand{\motivationtitlefont}{\normalsize\bfseries}
\renewcommand{\motivationbodyfont}{\footnotesize}
\renewcommand{\motivationpadding}{6pt}
\renewcommand{\motivationgap}{5pt}
\renewcommand{\motivationtextgap}{2pt}
\centering
\begin{minipage}{0.94\textwidth}
  \motivationcard{Regular path queries}
    {Describe graph paths using regular expressions over edge labels.}{UHblack!3}
  \motivationcard{Datalog\textsuperscript{$\circ$}}
    {Combine recursion with aggregation over semirings, e.g., shortest paths using min and addition.}{UHblack!3}
  \motivationcard{Rel}
    {Relational AI's language: recursive by design}{UHblack!3}
\end{minipage}
\par
\begin{tikzpicture}
  \draw[-latex,UHred,line width=1.5pt] (0,0.45) -- (0,0);
\end{tikzpicture}
\par\vspace{3pt}
\begin{tikzpicture}
  \node[fill=UHred!6,draw=UHred!40,rounded corners=4pt,
    text width=\dimexpr\textwidth-0.8cm-1pt\relax,
    inner xsep=0.4cm,inner ysep=9pt,align=center] {%
    {\color{red1}\large\bfseries Recursion within SQL\par}
    \vspace{3pt}
    {\small Express these computations in SQL and optimize their execution.\par}
  };
\end{tikzpicture}
% Rel: A Programming Language for Relational Data -- Aref et al., SIGMOD 2025.
\end{frame}

\begin{frame}{Why optimize recursive SQL?}
\vspace{0.6em}
{\small\color{UHblack!75}A recursive step repeats. Work independent of the recursion need not.\par}
\vspace{0.8em}
{\small\textbf{Example:} Find routes using only approved airlines.\par}
\vspace{0.7em}
\begin{columns}[T,onlytextwidth]
  \begin{column}{0.485\textwidth}
    \motivationcard{Repeated work}
      {At each iteration:\par
       build eligible flights;\par
       extend the current routes.}{UHblack!3}
  \end{column}
  \begin{column}{0.485\textwidth}
    \motivationcard{Loop hoisting}
      {Once: build eligible flights.\par
       Each iteration: extend routes\par
       using the prepared relation.}{UHred!6}
  \end{column}
\end{columns}
\vspace{0.5em}
\motivationcard{When is loop hoisting correct and beneficial?}
  {Identify invariant work, preserve SQL semantics, and measure the performance impact.}{UHblack!3}
% Illustrative execution strategies, not a claim about every SQL optimizer.
% Recursive evaluation reference: PostgreSQL documentation, section 7.8.2.
% https://www.postgresql.org/docs/current/queries-with.html
\end{frame}

\begin{frame}{Recursion in SQL}
\vspace{0.8em}
\renewcommand{\motivationpadding}{5pt}
\renewcommand{\motivationgap}{5pt}
\renewcommand{\motivationtitlefont}{\normalsize\bfseries}
\motivationcard{Recursive queries through recursive Common Table Expressions (CTEs)}
  {Part of the SQL standard since \textbf{SQL:1999}.\par}{UHblack!3}
\vspace{0.5em}
\centering
\motivationcard{Transitive closure}
  {{\ttfamily\small
  \begin{tabular}{@{}l@{}}
  \lstinline[style=slidesql]|WITH RECURSIVE tc(src, dst) AS (|\\
  \quad \lstinline[style=slidesql]|-- Anchor: direct edges|\\
  \quad \lstinline[style=slidesql]|SELECT src, dst FROM edge|\\
  \quad \lstinline[style=slidesql]|UNION|\\
  \quad \lstinline[style=slidesql]|-- Recursion: extend by one edge|\\
  \quad \lstinline[style=slidesql]|SELECT tc.src, e.dst|\\
  \quad \lstinline[style=slidesql]|FROM tc JOIN edge e ON tc.dst = e.src|\\
  \lstinline[style=slidesql]|)|\\
  \lstinline[style=slidesql]|SELECT src, dst FROM tc;|
  \end{tabular}}\par
  \vspace{4pt}
  {}}{UHred!6}

% \begin{tikzpicture}[
%   stage/.style={draw=UHblack!9,fill=UHblack!3,rounded corners=4pt,
%     text width=3.5cm,minimum height=1.3cm,inner sep=6pt,align=center},
%   flow/.style={-latex,UHred,line width=1.2pt}]
%   \node[stage] (anchor) at (2.0,0) {%
%     {\normalsize\bfseries Anchor term}\\[4pt]
%     {\footnotesize Produce the initial rows.}};
%   \node[stage,fill=UHred!6] (recursive) at (6.9,0) {%
%     {\normalsize\bfseries Recursive term}\\[4pt]
%     {\footnotesize Extend the latest rows.}};
%   \node[stage] (delta) at (11.8,0) {%
%     {\normalsize\bfseries Next working table}\\[4pt]
%     {\footnotesize Collect the next batch.}};
%   \draw[flow] (anchor.east) -- (recursive.west);
%   \draw[flow] (recursive.east) -- (delta.west);
%   \draw[flow] (delta.south) -- ++(0,-0.4)
%     -| node[pos=0.25,below,font=\footnotesize,text=UHblack!75]
%     {Repeat while nonempty} (recursive.south);
% \end{tikzpicture}
% \par\vspace{0.4em}
% \motivationcard{Semi-naive evaluation: process only the latest delta}
%   {Each iteration uses the previous round's new rows, rather than the full accumulated result.\par
%    {\footnotesize With \texttt{UNION}, discard duplicates and previously found rows; \texttt{UNION ALL} retains duplicates.}}{UHred!6}
% {\par\raggedright\tiny\color{UHblack!60}
% Sources: \href{https://www.postgresql.org/docs/current/queries-with.html\#QUERIES-WITH-RECURSIVE}{PostgreSQL, section 7.8.2} (linear recursion);
% \href{https://zedware.github.io/homepage/db/sql99-known-as-sql3.pdf}{SQL:1999, formerly known as SQL3}.\par}
\end{frame}

\section{Section 1: text}
\begin{frame}
\frametitle{My first frametitle}
Welcome to this DSI Latex Beamer template!
\vspace{2cm}
\\You can adjust the view of the section titles to repeating the table of contents with the current section highlighted. For this, follow the steps in line 316 of beamerthemeDSI.sty.
\end{frame}

\begin{frame}
You can also leave out a frametitle. The primary colors that you can use are
\begin{itemize}
\item \textcolor{UHred}{UHred}
\item \textcolor{UHblack}{UHblack}
\item \colorbox{UHblack}{\textcolor{UHwhite}{UHwhite}}
\end{itemize}

The secondary colors are 
\begin{itemize}
\item \textcolor{red1}{red1}
\item \textcolor{red2}{red2}
\item \textcolor{red3}{red3}
\item \textcolor{red4}{red4}
\end{itemize}
\end{frame}

\begin{frame}{Itemization}
\begin{itemize}
\item Item 1
\item Item 2
\begin{itemize}
\item Item 2.1
\item Item 2.2
\end{itemize}
\item Item 3
\end{itemize}
\end{frame}

\begin{frame}{Enumeration}
\begin{enumerate}
\item Item 1
\item Item 2
\begin{enumerate}
\item Item 2.1
\item Item 2.2
\end{enumerate}
\item Item 3
\end{enumerate}
\end{frame}

\section{Section 2: Blocks}

\subsection{Normal block}
\begin{frame}
Here is a normal block:

\begin{block}{Block title}
block body
\end{block}
\end{frame}

\subsection{Example block}
\begin{frame}
Here is an example block:

\begin{exampleblock}{Block title}
block body
\end{exampleblock}
\end{frame}


\subsection{Alert block}
\begin{frame}
Here is an alertblock:

\begin{alertblock}{Block title}
block body
\end{alertblock}
\end{frame}
\end{document}